\documentclass[12pt,letterpaper]{article}

\usepackage[utf8]{inputenc}
\usepackage[spanish]{babel}
\usepackage[margin=2.5cm]{geometry}
\usepackage{times}
\usepackage{amsmath}
\usepackage{listings}
\usepackage{xcolor}
\usepackage{booktabs}
\usepackage{array}
\usepackage{hyperref}
\usepackage{titlesec}
\usepackage{setspace}

\definecolor{codebg}{RGB}{245,245,245}
\definecolor{codekeyword}{RGB}{0,0,150}
\definecolor{codecomment}{RGB}{100,100,100}

\lstdefinestyle{cstyle}{
    backgroundcolor = \color{codebg},
    commentstyle = \color{codecomment},
    keywordstyle = \color{codekeyword}\bfseries,
    basicstyle = \ttfamily\small,
    breaklines = true,
    numbers = left,
    numberstyle = \tiny,
    frame = single,
    language = C
}
\lstset{style = cstyle}

\titleformat{\section}{\large\bfseries}{\thesection}{1em}{}
\titleformat{\subsection}{\normalsize\bfseries}{\thesubsection}{1em}{}

\title{Práctica \#1: Documento Técnico}
\author{Joanfer Hidalgo Chaves}
\date{\today}

\begin{document}
\shorthandoff{"}

\begin{titlepage}
\centering
\vspace*{1.5cm}

\Large{\textbf{Práctica \#1: Documento Técnico}}\\[6pt]
\large{Fundamentos de Eficiencia Computacional y Tipos de Datos Abstractos}\\[6pt]
Estructura Lineal: Arreglos Dinámicos\\

\vspace{1.5cm}

\textbf{Estudiante:}\\
Joanfer Hidalgo Chaves

\vspace{1.5cm}

\begin{spacing}{1.8}
Escuela de Informática y Computación\\
Universidad Nacional de Costa Rica\\
EIF207, Estructuras de Datos\\[10pt]
\textbf{Docente:} \\
Msc. Johnny Villalobos Murillo\\[1.5cm]
II Ciclo 2026\\
\today
\end{spacing}

\vfill
\end{titlepage}

\newpage
\tableofcontents

\begin{center}
\begin{minipage}{1\textwidth}
\centering
\small
\textbf{Recursos} \\[0.3em]
Este documento también forma parte del repositorio del proyecto, junto con el código fuente completo y su documentación generada con Doxygen:
\vspace{0.3em}

\begin{tabular}{ll}
Repositorio de GitHub: & \url{https://github.com/Cetrei/estructuras_practica1} \\
Documentación en línea (GitHub Pages): & \url{https://Cetrei.github.io/estructuras_practica1/} \\
\end{tabular}
\end{minipage}
\end{center}
\newpage

\section{Descripción de las Estructuras Utilizadas}

El programa define cuatro estructuras (\texttt{struct} y \texttt{enum}) para representar la información del ejercicio.
Todos los tamaños de cada atributo son fijos para poder leer y escribir los registros como se solicita sin generar errores.

\subsection{Estructura Estudiante}
Representa un registro del archivo \texttt{estudiantes.dat}.

\begin{lstlisting}
struct Estudiante:
    carnet[15]
    nombre[40]
    apellidos[40]
    carrera[40]
    nivel
\end{lstlisting}

\subsection{Estructura Historial}
Representa un registro del archivo \texttt{historial.dat}.

\begin{lstlisting}
struct Historial:
    carnet[15]
    codigoCurso[10]
    nombreCurso[50]
    ciclo
    anio
    nota
\end{lstlisting}

\subsection{Estructura Estadisticas}
Agrupa los resultados solicitados respecto a los datos almacenados en el historial.

\begin{lstlisting}
struct Estadisticas:
    notaMenor
    notaMayor
    notaPromedio
    carnetNotaMenor[15]
    carnetNotaMayor[15]
\end{lstlisting}

\subsection{Estructura CodigoResultado}
Representa el estado de salida de una operación que puede fallar de más de una forma.

\begin{lstlisting}
enum CodigoResultado:
    EXITO
    ARCHIVO_NO_ENCONTRADO
    MEMORIA_INSUFICIENTE
    NO_ENCONTRADO
\end{lstlisting}

\section{Algoritmo para Determinar la Cantidad de Registros}
\textbf{Función:} \texttt{contarRegistros(const char* nombreArchivo, int* cantidadRegistros)}

El algoritmo evita leer el contenido del archivo para saber cuántos registros contiene. 
En su lugar, se apoya en el tamaño del archivo en bytes:

\begin{enumerate}
    \item Abrir el archivo en modo de lectura (\texttt{"rb"}).
    \item Si el archivo no existe o no se pudo abrir, retornar \texttt{ARCHIVO\_NO\_ENCONTRADO}
    \item Mover el cursor de lectura al final del archivo
    \item Obtener la posición del cursor (el tamaño total en bytes)
    \item Cerrar el archivo
    \item Dividir el tamaño en bytes entre \texttt{sizeof(Historial)} para obtener la cantidad de registros completos almacenados
\end{enumerate}

\section{Algoritmo para Cargar el Arreglo Dinámico}
\textbf{Función:} \texttt{cargarHistorial(const char* nombreArchivo, int cantidadRegistros, Historial** historiales)}

\begin{enumerate}
    \item Abrir el archivo en modo de lectura
    \item Si no se pudo abrir, retornar \texttt{ARCHIVO\_NO\_ENCONTRADO}
    \item Reservar memoria dinámica para \texttt{cantidadRegistros} elementos de tipo \texttt{Historial}
    \item Si la reserva de memoria falla, cerrar el archivo y retornar \texttt{MEMORIA\_INSUFICIENTE}
    \item Leer todos los registros del archivo hacia el arreglo reservado indicando el tamaño de cada elemento y la cantidad total a leer.
    \item Cerrar el archivo y retornar \texttt{EXITO}
\end{enumerate}

La carga se realiza en una sola lectura en vez de una por registro ya que recibe la cantidad de elementos a leer y los copia
seguido en el arreglo aprovechando que el struct en memoria y archivo comparten el mismo tamaño por campo.

\section{Algoritmo para Calcular Nota Menor, Nota Mayor y Nota Promedio}
\textbf{Función:} \texttt{calcularEstadisticas(const Historial historiales[], int cantidadRegistros)}

En vez de recorrer el arreglo tres veces (una para el mínimo, otra para el máximo y otra para el promedio), 
el algoritmo realiza un único recorrido que actualiza los tres resultados simultáneamente.

\begin{enumerate}
    \item Inicializar \texttt{notaMenor} y \texttt{notaMayor} con la nota del primer registro, y sus carnets asociados con el carnet del primer registro
    \item Para cada registro del arreglo:
    \begin{enumerate}
        \item Sumar su nota a un contador \texttt{sumaNotas}.
        \item Si su nota es menor que \texttt{notaMenor}, actualizar \texttt{notaMenor} y \texttt{carnetNotaMenor}
        \item Si su nota es mayor que \texttt{notaMayor}, actualizar \texttt{notaMayor} y \texttt{carnetNotaMayor}
    \end{enumerate}
    \item Calcular \texttt{notaPromedio} dividiendo \texttt{sumaNotas} entre la cantidad de registros
\end{enumerate}

\section{Algoritmo para Localizar al Estudiante de Nota Menor y Nota Mayor}
\textbf{Función:} \texttt{buscarEstudiante(const char* nombreArchivo, const char* carnet, Estudiante* estudianteEncontrado)}

\begin{enumerate}
    \item Conseguir el carnet de la nota menor y el carnet de la nota mayor con \texttt{calcularEstadisticas()}
    \item Iniciar la busqueda para el carnet de la nota menor o mayor
    \item Abrir el archivo de estudiantes en modo de lectura.
    \item Si no se pudo abrir, retornar \texttt{ARCHIVO\_NO\_ENCONTRADO}
    \item Leer registros uno a uno sin cargar el archivo completo en memoria
    \item Comparar el carnet de cada registro leído contra el carnet buscado
    \item Si coinciden, copiar el registro encontrado en el parámetro de salida, cerrar el archivo y retornar \texttt{EXITO}
    \item Si se llega al final del archivo sin encontrar coincidencia, cerrar el archivo y retornar \texttt{NO\_ENCONTRADO}
\end{enumerate}

Se debe llamar esta funcion dos veces; una con el carnet de la nota menor y otra con el carnet de la nota mayor.

\section{Análisis del Costo Computacional}
\subsection{Supuestos}

Para el análisis se consideran las siguientes variables:

\begin{itemize}
    \item $n$ = cantidad de registros almacenados en \texttt{historial.dat}.
    \item $m$ = cantidad de registros almacenados en \texttt{estudiantes.dat}.
\end{itemize}

Para cada función se presenta el razonamiento que deduce su costo, distinguiendo:

\begin{itemize}
    \item \textbf{Complejidad algorítmica}: la cota asintótica reducida
    \item \textbf{Complejidad real}: el conteo efectivo de operaciones, sin simplificar, incluyendo constantes y pasos fijos
\end{itemize}

\subsection{contarRegistros()}

La función realiza una cantidad fija de operaciones sin importar el tamaño del archivo: 
\begin{itemize}
    \item \texttt{fopen}: abrir el archivo.
    \item \texttt{fseek}: mover el cursor al final del archivo.
    \item \texttt{ftell}: leer la posición del cursor.
    \item \texttt{fclose}: cerrar el archivo.
\end{itemize}
Ninguna de estas operaciones depende de $n$, porque \texttt{fseek} y \texttt{ftell} consultan el tamaño del archivo, no lo recorren
\newline \newline
\textbf{Complejidad real:} $O(4)$. $4$ operaciones de E/S (\texttt{fopen}, \texttt{fseek}, \texttt{ftell}, \texttt{fclose})
\newline
\textbf{Complejidad algorítmica:} $O(1)$. Ya que el costo no depende de $n$

\subsection{cargarHistorial()}

La función abre el archivo, reserva memoria para \texttt{cantidadRegistros} elementos, y realiza una única llamada que transfiere los $n$ registros del archivo al arreglo. 
La lectura en bloque copia los $n$ registros de forma continua, por lo que internamente crece proporcionalmente a $n$: cada uno de los $n$ registros se copia una vez.
\\[0.6cm]
\textbf{Complejidad real:} $O(n + 3)$. $1$ apertura, $1$ reserva de memoria, $n$ registros transferidos, $1$ cierre
\newline
\textbf{Complejidad algorítmica:} $O(n)$. El costo dominante es la transferencia de los $n$ registros

\subsection{calcularEstadisticas()}

La función recorre el arreglo de $n$ historiales una única vez. En cada iteración realiza una cantidad fija de operaciones: 
\begin{itemize}
    \item sumar la nota al contador de suma
    \item comparar contra la nota menor
    \item comparar contra la nota mayor
    \item y en el peor caso copiar el carnet en ambas comparaciones
\end{itemize}
Como estas comparaciones y copias no dependen de $n$ individualmente (son de tamaño fijo, 15 caracteres), el costo por iteración es constante, y se repite $n$ veces.
\\[0.6cm]
\textbf{Complejidad real:} $O(n)$. Al ser una iteracion lineal sobre $n$ elementos
\newline
\textbf{Complejidad algorítmica:} $O(n)$

\subsection{buscarEstudiante()}

La función recorre el archivo de estudiantes registro por registro, comparando el carnet de cada uno contra el carnet buscado, hasta encontrar coincidencia o llegar al final. 
En el peor caso (el carnet buscado es el último registro, o no existe), se leen y comparan los $m$ registros del archivo.
\\[0.6cm]
\textbf{Complejidad real:} $O(m + 2)$. $1$ apertura, hasta $m$ lecturas (\texttt{fread}) en el peor caso (no se encuentra o el registro buscado es el último), $1$ cierre
\newline
\textbf{Complejidad algorítmica:} $O(m)$

Esta función se invoca dos veces desde el flujo principal; una por el carnet de la nota menor y otra por el de la nota mayor.
Por lo que el costo real de las dos búsquedas combinadas es $O(m+2) + O(m+2) = O(2m + 4)$

\subsection{Costo Total del Programa}

Sumando el costo real de cada operación ejecutada desde el flujo principal (usando el desglose real de cada función, sin simplificar), incluyendo las operaciones de costo constante que arman e imprimen el reporte y la liberación final de memoria:

\[
O(4) + O(n + 3) + O(n) + O(m + 2) + O(m + 2) + O(1) + O(1) + O(1)
\]

Agrupando términos semejantes:

\[
O(2n + 2m + 14)
\]
\newline
\textbf{Complejidad real:} $O(2n + 2m + 14)$
\newline
\textbf{Complejidad algorítmica:} $O(n + m)$. Al eliminar constantes y factores menores

\subsection{Tabla Resumen}

\begin{center}
\begin{tabular}{@{}lc@{}}
\toprule
\textbf{Función} & \textbf{Complejidad algorítmica} \\
\midrule
\texttt{contarRegistros()} & $O(4)$ = $O(1)$ \\
\texttt{cargarHistorial()} & $O(n + 3)$ = $O(n)$ \\
\texttt{calcularEstadisticas()} & $O(n)$ \\
\texttt{buscarEstudiante()} (nota menor) & $O(m + 2)$ = $O(m)$ \\
\texttt{buscarEstudiante()} (nota mayor) & $O(m + 2)$ = $O(m)$ \\
\texttt{construirReporte()} & $O(1)$ \\
\texttt{imprimirReporte()} & $O(1)$ \\
\midrule
\textbf{Programa completo} & $O(n + m)$ \\
\textbf{Programa completo (real)} & $O(2n + 2m + 14)$ \\
\bottomrule
\end{tabular}
\end{center}

\end{document}